\documentclass[10pt]{examdesign}
\usepackage{amsmath}
\usepackage{pifont}
\SectionFont{\large\sffamily}
\Fullpages
\ContinuousNumbering
%\ShortKey
%\NoKey
\setrandomseed{8282023}
\DefineAnswerWrapper{}{}
\NumberOfVersions{3}

\class{MAT 2000: Introduction to Statistics \\ Grace College}
\examname{Quiz 1.1-3}
\def\namedata{Name : \hrulefill}

\begin{document}
	
	\begin{matching}
		Match the title of the six-step method with the descriptions.
		\pair{1. Ask a research question.}{Confirm the question can be addressed by collecting data. Typical questions involve comparing groups, asking how two characteristics are related, or surveying people's opinions.}
		\pair{2. Design a study and collect data.}{Plan how to select the people or objects to be studied, decide how to gather the relevant data, and carry out the data collection in a careful, systematic manner.}
		\pair{3. Explore the data.}{Look for patterns that relate to the research question or raise additional questions.}
		\pair{4. Draw inferences beyond the data.}{Determine whether any findings in the data reflect a genuine tendency and estimate the strength of that tendency.}
		\pair{5. Formulate conclusions.}{Consider whether inferences made in step 4 can be generalized to processes or groups beyond those represented in the study. Identify whether a cause-and-effect conclusion is warranted.}
		\pair{6. Look back and ahead.}{Point out limitations of the study and suggest new studies to build on the findings.}
	\end{matching}
	
	\begin{matching}
		Match each term from section 1.1 with it's definition
		\pair{data}{the values measured or categories recorded on individual objects of interest}
		\pair{observational unit}{the individual object on which data are recorded}
		\pair{variable}{a characteristic being recorded about the observational units}
		\pair{quantitative variable}{has numeric values for which ordinary arithmetic (addition, multiplication) makes sense}
		\pair{categorical variable}{places observational units into categories}
		\pair{distribution}{the pattern of outcomes of a variable}
	\end{matching}
	
	\pagebreak
	
	\begin{matching}
		Match each term from section 1.2 with it's definition
		\pair{dotplot}{a graph with one dot representing the variable outcome for each observational unit}
		\pair{histogram}{a graph that separates the quantitative variable into bins and displays the number of observations in each bin}
		\pair{shape of a distribution}{described by the number of peaks and whether one side of the distribution tends to stretch out longer than the other}
		\pair{center of a distribution}{the middle or typical value}
		\pair{variability of a distribution}{the spread in the observations}
		\pair{Unusual observations}{observations that do not follow the overall pattern of the distribution}
	\end{matching}
	
	\begin{matching}
		Match each term from section 1.3 with it's definition
		\pair{random process}{a repeatable process where the result for an individual attempt is not known in advance but there is a predictable pattern to the results in the long run}
		\pair{simulation}{an artificial re-creation of a random process}
		\pair{probability of an outcome}{the long-run proportion of times an outcome from a random process occurs if the process is repeated indefinitely under identical conditions}
	\end{matching}
	
	
	
\end{document}
